2. Data and descriptive evidence



\subsection{Data sources}



All data come from official sources, downloaded on 14 August 2026.


\textit{Homicides.} Incident-level records of intentional homicides (one row per victim) come from the DINASED registry of the Ministry of the Interior, published through the Ecuadorian open-data portal (datosabiertos.gob.ec). The dataset contains 39,822 records for 2014--2025 plus 4,154 for January--June 2026 (updated 15 July 2026). Records were aggregated to the province--quarter level using the two-digit DPA (División Político-Administrativa) provincial code as the join key.


\textit{Population.} Denominators are INEC population projections (2024 revision, 2022 Census base; 1990--2035 series), used to compute homicide rates per 100,000 inhabitants. These are projections, not census counts, and therefore introduce a source of measurement error that is likely modest for provincial rates over this horizon.


\textit{Labor market.} Provincial unemployment, underemployment, and adequate employment rates are computed from ENEMDU (Encuesta Nacional de Empleo, Desempleo y Subempleo) microdata for 2024-II, the only quarter for which the project has processed provincial cross-sections, using expansion factors (\texttt{fexp}) and validated against the official INEC bulletin. Official tabulations do not provide provincial disaggregation, which is why microdata processing is required.


\textit{Geography.} Provincial boundaries (24 provinces) come from the CONALI/IGM territorial organization file for 2025 (EPSG:32717, reprojected to EPSG:4326).


\textit{External contrast.} The OECO (Observatorio Ecuatoriano de Crimen Organizado, PADF) semi-annual bulletin S1-2025 reports 4,619 intentional homicides for the first half of 2025, against 4,659 in the official dataset used here -- a difference of approximately 0.9\%, attributable to cut-off dates. This provides an independent validation of the primary source.



\subsection{Panel construction}



The analysis uses a balanced panel of 24 provinces $\times$ 34 quarters = 816 observations spanning 2018Q1--2026Q2. Homicide counts are aggregated by province and quarter of occurrence; rates per 100,000 are computed with the corresponding annual population projection. Labor indicators are provincial rates from the 2024-II cross-section -- time-invariant in the panel, which determines the identification strategy of Section 3. Galápagos reports no homicides in 2026; it is retained in the panel with zero counts for those quarters.



\subsection{The national series}



Figure 1 presents the national series of intentional homicides (annual counts, 2018--2025, and the first-half 2026 partial count). The series documents four phases: relative stability around 996 victims in 2018; a first acceleration beginning in 2021 (2,495 victims); an explosive rise through 2023 (8,248 victims); and persistently high levels in 2024--2025, culminating in 9,283 victims in 2025. The 2026 first semester registers 4,154 homicides, slightly below the 4,659 recorded in the first half of 2025. The national rate reached approximately 51 per 100,000 in 2025, among the highest in Latin America.


\begin{figure}[htbp]
  \centering
  \includegraphics[width=0.92\textwidth]{fig01_national_series.png}
  \caption{Intentional homicides in Ecuador, 2018--2026. Annual counts (2026 shown as first-semester partial, in red). Source: DINASED (Ministry of the Interior); authors' calculations.}
  \label{fig:serie}
\end{figure}



\subsection{Provincial concentration and geographic breadth}



The wave is heavily concentrated but not confined to one province. Guayas, home to the port of Guayaquil, accounts for 44.2\% of accumulated homicides over 2018--2025, yet re-estimating the series excluding Guayas shows the same upward direction: homicides outside Guayas rose 34.7\% in 2025 relative to 2024, and provincial maps show deterioration extending along the coast (Los Ríos, Manabí, El Oro, Esmeraldas) and beyond. The crisis is a national phenomenon with a strong epicenter, not a Guayas-specific anomaly.

\begin{figure}[htbp]
  \centering
  \includegraphics[width=0.62\textwidth]{fig02_map_2025.png}
  \caption{Provincial homicide rates per 100,000 inhabitants, 2025. Source: DINASED records and INEC population projections; CONALI/IGM boundaries.}
  \label{fig:mapa}
\end{figure}




\subsection{The labor market profile}



The ENEMDU 2024-II national profile is one of low open unemployment but widespread precariousness: unemployment 3.48\%, underemployment 20.97\%, and adequate employment 35.01\%. This is the canonical profile of a highly informal labor market, in which the opportunity cost of illegal activity is determined less by the unemployment rate than by the availability and quality of licit income (Fajnzylber et al., 2002a; Britto et al., 2022). Provincial variation in these indicators, together with 2025 provincial homicide rates, is the raw material for the correlations of Section 4.


Table~\ref{tab:descriptivos} reports the provincial distribution of the three labor indicators and of the 2025 homicide rate (N = 24 provinces). Provincial unemployment averages 3.05\% (SD 2.20, range 0.37--8.91), underemployment 19.77\% (SD 5.85, range 6.81--30.72), and adequate employment 29.58\% (SD 11.33, range 8.78--52.67); the 2025 provincial homicide rate averages 35.29 per 100,000 (SD 36.06, range 0--130.41). Only the national labor figures are verified against the official bulletin: unemployment 3.48\%, underemployment 20.97\%, and adequate employment 35.01\% (INEC, ENEMDU 2024-II).

\begin{table}[htbp]
  \centering
  \small
  \caption{Provincial distribution of labor indicators (ENEMDU 2024-II) and homicide rates (2025): descriptive statistics, N = 24 provinces.}
  \label{tab:descriptivos}
  \begin{tabular}{lcccc}
    \toprule
    Variable & Mean & SD & Min & Max \\
    \midrule
    Unemployment (\% of labor force, 2024-II) & 3.05 & 2.20 & 0.37 & 8.91 \\
    Underemployment (\%, 2024-II) & 19.77 & 5.85 & 6.81 & 30.72 \\
    Adequate employment (\%, 2024-II) & 29.58 & 11.33 & 8.78 & 52.67 \\
    Homicide rate 2025 (per 100,000) & 35.29 & 36.06 & 0.00 & 130.41 \\
    \bottomrule
  \end{tabular}
  \par\vspace{3pt}\noindent\small\textit{Note. }Source: ENEMDU 2024-II microdata (INEC) and DINASED homicide records, authors' calculations. Homicide rates computed with INEC population projections. Galápagos records zero homicides and is included with rate 0.
  
\end{table}



\subsection{Limitations of the data}



Four limitations are acknowledged. First, labor indicators are available for a single quarter, so the panel cannot support temporal labor variation or long-run elasticity estimates. Second, homicide records are administrative and "subject to variations" per official metadata; they exclude other crimes and unclassified violent deaths, and Galápagos reports zero homicides in 2026. Third, population denominators are projections. Fourth, the open-data portal's API blocks automated downloads, constraining replication.